\printedpage{99}
% head: A. Grothendieck : Dérivateurs, Ch. VII, Édition des Universités de Montpellier II et Paris VII
we will also have to prove the results corresponding to 1) 2) for $\cC_{\NN}$, i.e.

\begin{itemize}[label={}]
\item[1$'$)]
\[
\begin{tikzcd}
X \ar[r] \ar[d,"\wr"'] & \cC_{\NN} \\
\cC_{t_0} = \cC_0 &
\end{tikzcd}
\]
(the canonical inclusion of $\cC_0$, fiber of $\cC_{\NN}$ (fibered over $\mathbf{N}^o$) at $0 \in \mathbf{N}^o$) is a Hot-equivalence, and:

\item[2$'$)] $\cC_{\NN} \longrightarrow X \times X$ is proper and smooth.
\end{itemize}

With these four technical results, which I am going to prove shortly (and $1'$ and $2'$ come as no surprise to me), we win. Let us indeed consider the diagram
\[
\begin{tikzcd}
& X \ar[dl,"\alpha_1"'] \ar[d,"\alpha_0"] \ar[dr,"\alpha_2"] & \\
\Ch(X) \ar[dr,"p_1"'] & \cC_{\NN} \ar[l,"\pi_1"'] \ar[r,"\pi_2"] \ar[d,"p_0"] & \Ch_{\infty}(X) \ar[dl,"p_2"] \\
& X \times X. &
\end{tikzcd}
\]

By 1) and $1'$) (concerning $\alpha_1$ and $\alpha_0$) and the similar result
\msp{page 159}
already known for $\alpha_2$, these three functors are Hot-equivalences. So the same holds for $\pi_1$, $\pi_2$. Since $p_1$, $p_0$, $p_2$ are Hot-fibrations (being proper and smooth), it follows that $\pi_1$ and $\pi_2$ induce Hot-equivalences on the fibers
\[
\Ch(X;a,b) \quad \xleftarrow{\ \pi_1(a,b)\ } \quad \cC_{\NN}(a,b) \quad \xrightarrow{\ \pi_2(a,b)\ } \quad \Ch_{\infty}(X;a,b).
\]
I had checked this (without difficulty) already for $\pi_2$ and the $\pi_2(a,b)$, which had made it possible, by inverting $\pi_2$ in Hot, to define \textit{in} Hot
\[
\left\{
\begin{aligned}
\Ch_{\infty}(X) &\quad\longrightarrow\quad \Ch(X) \\
\Ch_{\infty}(X;a,b) &\quad\longrightarrow\quad \Ch(X;a,b),
\end{aligned}
\right.
\]
and to say that these are isomorphisms means precisely that $\pi_1$ and the $\pi_1(a,b)$ are isomorphisms.

I do not think that $\cC_{\NN}$ is in itself of any other interest, than to make it possible to compare % [sic] interêt -> intérêt, missing acute accent
the $X \times X$-categories $\Ch_{\infty}(X)$ and $\Ch(X)$, with the help of the preceding diagram. On the other hand, $\Ch_{\infty}(X)$ and $\Ch(X)$ appear to me as the two essential ingredients of the formalism of % [sic] m'aparaissent -> m'apparaissent, spelling
paths in Cat.
% foot: Version 22 décembre 2001

\printedpage{100}
% head: A. Grothendieck : Dérivateurs, Ch. VII, Édition des Universités de Montpellier II et Paris VII
\msp{page 160}
I am going to prove 1, 2, $1'$, $2'$.

1) Let us consider the diagram
\[
\begin{tikzcd}
X = (X_{\Ty})_e \ar[rr,hook,"i"] \ar[dr,dotted] & & X_{\Ty} = X \times \Ty \ar[rr,"j"] \ar[dr] & & \Ch(X) \ar[dl] \\
& e \ar[rr,dotted] & & \Ty\ , &
\end{tikzcd}
\]
where $X_{\Ty}$ and $\Ch(X)$ are considered as Cat-fibered categories over $\Ty$, and the functor $j$ is deduced from the canonical functors
\[
X \longrightarrow \Ch(X)_{\tau} \overset{\text{def}}{=} \Ch_{\tau}(X),
\]
associating to $x \in \Ob X$ the path $p_{\tau}^{*}(x)$, where $p_{\tau} : \tau \longrightarrow e$ is the projection from $\tau$ to the unit element (and final element) $e = 1_{\Ty}$ of $\Ty$, and $\Ch_{e}(X)$ is identified with $X$. We have to check the

\textbf{Lemma.} \textit{For every $\tau \in \Ty$, the preceding functor is a} Hot-\textit{equivalence (and even a homotopism).}

But we know that for two categories $\cX$, $\cC$ smooth over another, $T$, every $T$-morphism from $\cX$ to $\cC$ that induces Hot-equivalences on the fibers, is a Hot-equivalence. Thus $j$ is
\msp{page 161}
is a Hot-equivalence. On the other hand, $\Ty$, having a final object, is aspheric, hence % [sic] est ... est -> est, word repeated across the manuscript page break
$X \times \Ty \xrightarrow{\mathrm{pr}_1} X$ is a Hot-equivalence, whence it follows at once that $i$ is a Hot-equivalence. So likewise $ji$, now $ji = \alpha_1$, q.e.d.

2) Let us prove e.g. the smoothness of $\Ch(X) \longrightarrow X \times X$, the property will follow from it by applying this result to $X^{\mathrm{opp}}$. So let
\[
(a',b') \ \longrightarrow\ (a,b), \qquad \text{i.e. } a' \xrightarrow{u} a,\ \ b' \xrightarrow{v} b,
\]
be an arrow in $X \times X$, and $c$ an object of $\Ch(X)$ over $(a,b)$:
\[
\begin{tikzcd}[column sep=large]
a' \ar[r,dotted,"c'","\scriptstyle\rightarrow"'] \ar[d,"u"'] & b' \ar[d,"v"] \\
a \ar[r,"c","\scriptstyle\rightarrow"'] & b.
\end{tikzcd}
\]

We look at the category $\cC(c;u,v)$ of pairs $(c',w)$ that lift $(c,(u,v))$, where therefore $c'$ is a path (of unspecified type) from $a'$ to $b'$, and
\[
w : c' \longrightarrow c
\]
an arrow in $\Ch(X)$ which on the source objects induces $u$, and on the target objects induces $v$ (thus $c'$ is a path from $a'$ to $b'$). We have to
% foot: Version 22 décembre 2001
\printedpage{101}
% head: A. Grothendieck : Dérivateurs, Ch. VII, Édition des Universités de Montpellier II et Paris VII
\msp{page 162}
prove that this category $\cC(c;u,v)$ is aspheric. For this, we look at the composite path
\[
\tilde{c} = \mathrm{ch}(u)\ c\ \mathrm{ch}(v)^{o},
\]
where $\mathrm{ch}(u)$ is the (direct) path from $a'$ to $a$ defined by $u : a' \longrightarrow a$, and $\mathrm{ch}(v)^{o}$ therefore denotes the (retrograde) path from $b$ to $b'$ defined by $v : b' \longrightarrow b$. If $c$ is of type $\tau$, $\tilde{c}$ is therefore of type $\Delta_{+}\tau\Delta_{-}$. We are going to compare $\cC(c;u,v)$ and the category $\Ch(X)/\tilde{c}$, with the help of two functors
\[
\begin{tikzcd}
\cC(c;u,v) \ar[r,shift left,"\varphi"] & \Ch(X;a',b')/\tilde{c}\ , \ar[l,shift left,"\psi"]
\end{tikzcd}
\]
which we will prove define arrows inverse to one another in Hot. Since $\Ch(X)/\tilde{c}$ is contractible (having a final object), hence aspheric, the same holds for $\cC(c;u,v)$, and we will win.

We define $\varphi$ by $\boxed{\varphi(c',w') = (\tilde{c}',\tilde{w}'),}$ where ($c'$ being of type $\tau'$) $\tilde{c}'$ is the path of type $\Delta_{+}\tau'\Delta_{-}$ obtained from $c'$ as inverse image under the degeneracy $\Delta_{+}\tau'\Delta_{-} \longrightarrow \tau'$, i.e. we extend $c'$ by a ``$\Delta_{+}$ step'' $a' \xrightarrow{\mathrm{id}} a'$ on the left and by a ``$\Delta_{-}$ step'' $b' \xleftarrow{\mathrm{id}} b'$ on the right, cf. figure.
\[
\overbrace{\underbrace{\begin{tikzcd}[ampersand replacement=\&]
a' \ar[r,"\mathrm{id}"] \ar[d,"\mathrm{id}"'] \ar[dr,phantom,"\ast"] \& a' \ar[rr,dotted,no head,"c'"] \ar[d,"u"'] \& {} \ar[d,dotted,"w'"] \& b' \ar[d,"v"'] \ar[dr,phantom,"\ast"] \& b' \ar[l,"\mathrm{id}"'] \ar[d,"\mathrm{id}"] \\
a' \ar[r,"u"] \& a' \ar[rr,no head,"c'"'] \& {} \& b' \& b' \ar[l,"v"'] % [sic] a' ... b' with label c' in the bottom row -> a ... b with label c, the bottom row of the figure is the path c~ = ch(u) c ch(v)^o
\end{tikzcd}}_{\tilde{c}}}^{\tilde{c}'}
\]

\msp{page 163}
We can say that
\[
\tilde{c}' = a'_{\Delta_{+}} \circ c' \circ b'_{\Delta_{-}},
\]
where in general, if $x \in X$ and $\tau \in \mathrm{Ty}$, $x_{\tau}$ denotes the ``constant path'' of value $x$, type $\tau$. On the other hand we have obvious morphisms
\[
a'_{\Delta_{+}} \longrightarrow \mathrm{ch}(u), \qquad b'_{\Delta_{-}} \longrightarrow \mathrm{ch}(v)^{\mathrm{opp}},
\]
as shown [?] by the vertical arrows $(\mathrm{id}_{a'},u)$ and $(v,\mathrm{id}_{b'})$ in the two squares marked $\ast$ in the diagram, arrows which lie one in $\Ch_{\Delta_{+}}(X)$, the other in $\Ch_{\Delta_{-}}(X)$. This said, we set
\[
\tilde{w}' : \tilde{c}' \longrightarrow \tilde{c}, \qquad \tilde{w} = \lambda(u) \ast w' \ast \mu(v), % [sic] \tilde{w} -> \tilde{w}', the morphism just named is \tilde{w}'
\]
where
\[
\left\{
\begin{aligned}
\lambda(u) :\ \ a'_{\Delta_{+}} &\quad\longrightarrow\quad \mathrm{ch}(u) \\
\mu(v) :\ \ b'_{\Delta_{-}} &\quad\longrightarrow\quad \mathrm{ch}(v)^{\mathrm{opp}}
\end{aligned}
\right.
\]
% foot: Version 22 décembre 2001
\printedpage{102}
% head: A. Grothendieck : Dérivateurs, Ch. VII, Édition des Universités de Montpellier II et Paris VII
(cf. figure). The pair
\[
(\tilde{c}',\tilde{w}') \in \Ob \Ch(X)/\tilde{c}
\]
thus defined depends functorially on $(c',w')$ varying in $\cC(c;u,v)$: If $(c'_1,w'_1)$ is a second object of this category, an arrow from the first to the second is an arrow
\[
c' \xrightarrow{\alpha} c'_1
\]
in $\Ch(X;a',b')$, such that we have commutativity in
\[
\begin{tikzcd}
c' \ar[rr,"\alpha"] \ar[dr,"w'"'] & & c'_1 \ar[dl,"w'_1"] \\
& c, &
\end{tikzcd}
\qquad\qquad w' = w'_1\alpha.
\]

Such an arrow $\alpha$ defines
\[
\tilde{\alpha} = a'_{\Delta_{+}} \ast \alpha \ast b'_{\Delta_{-}} : \tilde{c}' \longrightarrow \tilde{c}'_1
\]
in $\Ch(X;a',b')$, and we check
\msp{page 164}
at once that we have commutativity in
\[
\begin{tikzcd}
\tilde{c}' \ar[rr,"\tilde{\alpha}"] \ar[dr,"\tilde{w}'"'] & & \tilde{c}'_1 \ar[dl,"\tilde{w}'_1"] \\
& \tilde{c}, &
\end{tikzcd}
\qquad\qquad \text{i.e. } \tilde{w}' = \tilde{w}'_1\tilde{\alpha},
\]
i.e. that $\tilde{\alpha}$ is indeed an arrow in $\Ch(X;a',b')/\tilde{c}$. Likewise, that $\alpha \longmapsto \tilde{\alpha}$ defines a functorial dependence, i.e. $(\beta\alpha)\tilde{} = \tilde{\beta}\tilde{\alpha}$, is obvious from the definition of $\tilde{\alpha}$. So, we have indeed defined % Reading resolved against the supplied PDF; see corrections.csv.
\[
\cC(c;u,v) \ \xrightarrow{\ \varphi\ }\ \Ch(X;a',b')/\tilde{c},
\]
let us now define
\[
\cC(c;u,v) \ \xleftarrow{\ \psi\ }\ \Ch(X;a',b')/\tilde{c}.
\]
Let $(c'',w'')$ be an object of $\Ch(X;a',b')/\tilde{c}$,
\[
\overbrace{\underbrace{\begin{tikzcd}[ampersand replacement=\&]
a'_0 = a' \ar[r] \ar[d,"\mathrm{id}"'] \& a'_1 \ar[rr,no head,"\rightarrow"] \ar[d,"w''_1"'] \& {} \ar[d,dotted,"w''"] \& a'_n \ar[d,"w''_n"'] \& b' \ar[l] \ar[d,equal] \\
a' \ar[r,"u"] \& a \ar[rr,no head] \& {} \& b \& b', \ar[l,"v"']
\end{tikzcd}}_{c}}^{c''}
\]
% foot: Version 22 décembre 2001

\printedpage{103}
% head: A. Grothendieck : Dérivateurs, Ch. VII, Édition des Universités de Montpellier II et Paris VII
thus
\[
w'' : c'' \longrightarrow \tilde{c}, \qquad\quad \text{arrow in } \Ch(X;a',b'),
\]
let us note that we have an obvious retraction arrow
\[
\tilde{c} \xrightarrow{\rho_c} c \qquad\quad \text{associated with } \Delta_{+}\tau\Delta_{-} \longrightarrow \tau \text{ (retraction),}
\]
inducing $u$ and $v$ on the source and the target:
\[
\begin{tikzcd}
a' \ar[r,"u"] \ar[d,"u"] & a \ar[rrr,no head,"c"] \ar[d,"\mathrm{id}"] & {} \ar[d,"\mathrm{id}"] & {} \ar[d,"\mathrm{id}"] & b \ar[d,"\mathrm{id}"] & b' \ar[l,"v"'] \ar[d,"v"] & \Big\}\ \tilde{c} & \in \Ch_{\Delta_{+}\tau\Delta_{-}}(a',b') \\
a \ar[r,"\mathrm{id}"] & a \ar[rrr,no head,"c"] & {} & {} & b & b' \ar[l,"v"'] & \Big\}\ \delta^{*}_{\tau}(c) & \in \Ch_{\Delta_{+}\tau\Delta_{-}}(a,b) \\
& a \ar[rrr,no head,"c"] & & & b & & \Big\}\ c & \in \Ch_{\tau}(a,b).
\end{tikzcd}
\]

\msp{page 165}
So we define an object
\[
(c',w') = \psi(c'',w'')\ \in \Ob \cC(c;u,v)
\]
by
\[
c' = c'', \qquad\quad w' = \rho_c w'',
\]
and this is functorial in $(c'',w'')$.

The composite $\varphi\psi$ being an endomorphism of the contractible category $\Ch(X;a',b')/\tilde{c}$, is homotopic to the identity. We have to check that the same is true of $\psi\varphi$, an endomorphism of $\cC(c;u,v)$. This functor associates to $(c',w')$, $(\tilde{c}',\overline{w'})$, where $\tilde{c}'$ is defined as above, $\tilde{c}' = a'_{\Delta_{+}} \circ c' \circ b'_{\Delta_{-}}$, and where $\overline{w'}$ is the composite
\[
\underbrace{\tilde{c}'}_{= \delta^{\prime *}_{\tau'}(c')}\ \ \underbrace{\xrightarrow{\ \rho_{c'}\ }}_{\substack{\text{canonical}\\ \text{arrow}}}\ \ c' \xrightarrow{w'} c.
\]

So we have
\[
(\psi\varphi)(c',w') = (\tilde{c}',\overline{w'}) \xrightarrow{\rho_{c'}} (c',w'),
\]
this is indeed a morphism in $\cC(c;u,v)$, i.e. we have commutativity in
\[
\begin{tikzcd}
\tilde{c}' \ar[rr,"\rho_{c'}"] \ar[dr,"\overline{w'}"'] & & c' \ar[dl,"w'"] \\
& c &
\end{tikzcd}
\qquad\qquad \text{(arrow in } \Ch(X;a',b')\text{).}
\]
% foot: Version 22 décembre 2001